% Teil II: durchgerechneter Weg, Teile a--d.

\subsection*{Gegeben und gesucht}
Seil der L"ange $L$, Dichte $\rho$, Querschnitt $A$, Zugkraft $F$; kleine
Auslenkungen. Am Ende $x=0$ erzeugt das M"adchen das Wellenpaket
$y(0,t)=y_0\sin^2(\pi\,t/\mathrm{s})$ f"ur $0\le t\le\qty{1}{s}$, sonst $0$
($\mathrm{s}$ symbolisch f"ur \qty{1}{s}). Gesucht: (a) $c$, (b) Skizze und
Maximum bei $t=1$~s, (c) $y(x,t)$ der vorwärtslaufenden Störung, (d,
Bonus) Überlagerung mit entgegengesetzt ausgelenktem anderen Ende.

\subsection*{Ansatz und Physik dahinter}
Aus der Dynamik: je gr"o"ser die Rückstellkraft $F$, desto schneller, je
gr"o"ser die Tr"agheit $\mu=\rho A$, desto langsamer --- daraus direkt
$c=\sqrt{F/\mu}$. Das Paket wandert unverändert als Funktion von $t-x/c$
nach rechts. Für das Maximum setzt man die ausgesandte Störung in Bewegung:
ein Punkt, der zur Zeit $t_s$ ausgesandt wurde, ist zur Zeit $t$ bei
$x=c\,(t-t_s)$.

\subsection*{Rechnung}
\textbf{a) Ausbreitungsgeschwindigkeit.}
\begin{equation}
  c=\sqrt{\frac{F}{\mu}}=\sqrt{\frac{F}{\rho A}}\;.
\end{equation}

\textbf{b) Maximum bei $t=\qty{1}{s}$.} Das Paket wird w"ahrend
$t_s\in[0,1~s]$ ausgesandt; die Vorderseite ($t_s=0$) ist zur Zeit $t$ bei
$x=ct$. Die Auslenkungsform $\sin^2(\pi t_s/\mathrm{s})$ hat ihr Maximum bei
$t_s=\qty{0.5}{s}$ (Scheitel des $\sin^2$). Das Maximum wurde also zur Zeit
$t_s=0.5$~s ausgesandt und bewegt sich mit $c$. Zum Zeitpunkt $t=1$~s:
\begin{equation}
  x_{\max}=c\,(t-t_s)=c\,(1-0.5)\,\mathrm{s}
  =\tfrac12\,c\cdot\mathrm{s}
  \;.
\end{equation}
Mit $c$ in \unit{\metre\per\second} und $\mathrm s=\qty{1}{s}$ ist das die
L"ange $x_{\max}=\tfrac12\,c\cdot\qty{1}{s}$ (in Metern, mit $c$ als
Zahlenwert).
Das Maximum liegt zur halben Paketl"ange: zu $t=\qty{1}{s}$ ist das Paket
gerade fertig ausgesandt, reicht von $x=0$ bis $x=c\cdot\qty{1}{s}$, und sein
Berg liegt in der Mitte. Skizze: vor Ort gebundener, glockenf"ormiger Berg von
H"ohe $y_0$ bei $x=c/2$, rechts und links abfallend auf $0$.

\textbf{c) Laufende Störung (Richtung Junge, $+x$).} Jeder Punkt läuft mit
Zeitverz"ogerung $x/c$:
\begin{equation}
  y(x,t)=\begin{cases}
    y_0\,\sin^2\!\bigl(\pi\,(t-x/c)/\mathrm{s}\bigr)
      & \text{f"ur } 0\le t-x/c\le\qty{1}{s},\\
    0 & \text{sonst}.
  \end{cases}
\end{equation}
Die Bedingung $0\le t-x/c\le1$~s bedeutet: der Punkt $(x,t)$ ist "uberhaupt
erst erreicht, wenn die Vorderseite ($t_s=0$, also $x\le ct$) vorbeikam, und
schon wieder vorbei, nachdem die Rückseite ($t_s=1$~s, also $x\ge c(t-1~s)$)
durch ist. Es ist exakt die Funktion $y(0,t)$ mit Argument $t-x/c$.

\textbf{d, Bonus) Zwei gegenl"aufige Pakete.} Das M"adchen sendet nach
rechts $y_G(x,t)=y(0,t-x/c)$, der Junge nach links
$y_J(x,t)=-y(0,t-(L-x)/c)$. Beide haben r"aumliche L"ange $c\cdot\qty{1}{s}$.

Zum Zeitpunkt $t_1=L/2c$: die beiden Vorderkanten treffen sich gerade in der
Mitte $x=L/2$; es "uberlappen sich noch keine Störungen (die Pakete berühren
sich nur). Kurz danach schieben sie sich "ubereinander --- mit
\emph{entgegengesetztem Vorzeichen}.

Zum Zeitpunkt $t_2=t_1+\qty{0.5}{s}$: jedes Paket ist um weitere
$c\cdot\qty{0.5}{s}$ gerückt, und die beiden wären jetzt eine volle
Paketlänge verschoben --- sie \emph{decken sich exakt} und heben sich wegen
$y_J=-y_G$ \emph{vollst"andig auf}. Das Seil ist zu diesem Zeitpunkt
\emph{gerade} ($y\equiv0$), ein eindrücklicher Moment: zwei Maxima
verschwinden, wo sie sich treffen, die Energien der beiden ist trotzdem nur
überlagert (Superposition), die Pakete laufen ungestört weiter.

Zum Zeitpunkt $t_3=t_1+\qty{1}{s}$: die Pakete sind vollst"andig
\textbf{durchgelaufen} und getrennt --- das M"adchen-Paket ist jetzt auf der
rechten, das Jungen-Paket auf der linken Seite, jedes mit vollem $|y_0|$
(entgegengesetztes Vorzeichen). Das Seil zeigt zwei nach au"sen (nach Mitte
gegenl"aufig) versetzte Berge.

\begin{ergebnis}
(a) $c=\sqrt{F/(\rho A)}$. (b) Maximum bei $x=\frac12 c\cdot\qty{1}{s}$
(Berg in Paketmitte). (c) $y=y_0\sin^2(\pi(t-x/c)/\mathrm{s})$ f"ur
$0\le t-x/c\le\qty{1}{s}$, sonst $0$. (d) Bei $t_2$ exakte Überlagerung
$\rightarrow$ gerades Seil; bei $t_3$ getrennte, entgegengesetzt ausgelenkte
Pakete.
\end{ergebnis}

\begin{kontrolle}
\textbf{Einheiten:} $F/\mu=(\mathrm{kg\,m/s^2})/(\mathrm{kg/m})=\mathrm{m^2/s^2}$,
Wurzel $\rightarrow\unit{\metre\per\second}$. Das Argument $t-x/c$ von
$\sin^2$ ist dimensionslos. \checkmark

\textbf{Grenzfall:} F"ur $F\rightarrow0$ erg"abe $c\rightarrow0$ (loses Seil,
keine Welle), f"ur $\mu\rightarrow0$ (imagin"ar masseloses Seil) $c\rightarrow
\infty$. Beides physikalisch sinnvoll. \checkmark

\textbf{Querprobe (Anschaulich):} Ein Punkt auf dem Seil kann erst schwingen,
nachdem die Vorderseite ihn erreicht hat -- mit
$x_{\text{Reichweite}}=c\,t$ passt das zu Teil (b): zur Zeit $t=1$~s ist die
Reichweite gerade $c\cdot\qty{1}{s}$, das Maximum liegt bei halber Reichweite.
\checkmark
\end{kontrolle}

\begin{merksatz}
Eine gespannte Saite transportiert Transversalwellen mit
$c=\sqrt{F/\mu}$: R"uckstellkraft über Tr"agheit. Ein "uber ein Zeitintervall
$\Delta t$ erzeugtes Paket hat L"ange $c\,\Delta t$; zwei entgegengesetzte
Pakete löschen sich bei exakter Überlagerung aus und laufen danach ungestört
weiter (Superpositionsprinzip).
\end{merksatz}
